\section{The finite variational spectral step}
\label{sec:uncle_variational_spectral_step}

\begin{lemma}[A small positive eigenvalue from a trial vector]
    \label{lem:uncle_variational_positive_eigenvalue}
    \lean{MPSTensor.exists_pos_eigenvalue_uncleHamiltonianES_lt_of_rayleigh}
    \notready
    \uses{thm:uncle_kernel_local_conditions}
    Let $H'_P$ be the actual finite periodic uncle Hamiltonian. If
    $v\in(\ker H'_P)^\perp$ and
    \[
        \operatorname{Re}\langle H'_Pv,v\rangle<\varepsilon\|v\|^2,
    \]
    then $H'_P$ has a real eigenvalue $\mu$ with
    $0<\mu<\varepsilon$.
\end{lemma}
\begin{proof}
    The actual Hamiltonian is positive. Expand $v$ in an orthonormal
    eigenbasis, with nonnegative eigenvalues $\lambda_j$. A zero
    eigenvalue has zero coefficient because $v$ is orthogonal to the
    kernel. If every positive $\lambda_j$ were at least $\varepsilon$,
    the quadratic energy would be at least
    $\varepsilon\sum_j|v_j|^2$, contradicting the strict bound.
\end{proof}

The norm, energy, and ground-state overlap estimates must still construct
such a vector in the kernel complement. This lemma supplies only the final
finite spectral argument of
\cite[proof of the gaplessness theorem]{FernandezGonzalez2015Uncle}.
